\chapter{Bibliographical Study}
\markboth{Bibliographical Study}{Bibliographical Study}
\label{Bibliographical Study}

\section{UWB}

\subsection{Definitions} 
Ultra-Wideband (UWB) is a technology using radio frequency to exchange data.
Communication is not its only purpose; UWB can also be used to locate, which is only possible by using multiple modules. These algorithms and techniques will be covered in the next part (\ref{UWB_Algo_localisation}).
In the majority of cases, UWB is deployed for indoor purposes; otherwise, it can be used for outdoor purposes, but this is less present in scientific papers.


\subsection{Location algorithm}\label{UWB_Algo_localisation}
In order to use UWB as a location sensor, we have to use specific algorithms. In the following subsections, we will explain and show 2 different location algorithms: firstly, two-way ranging (TWR), and secondly, double-sided two-way ranging (DS-TWR). Finally, we will describe TDMA, the most advanced algorithm.

\subsubsection{Two way ranging (TWR)}
The Two-Way Ranging (TWR) algorithm is the easiest and the fastest to implement in a UWB system.
In order to explain it, we will use 2 modules called $UWB_1$ and $UWB_2$.

\begin{figure}[!ht]
\centering
\resizebox{0.7\textwidth}{!}{%
\begin{circuitikz}
\tikzstyle{every node}=[font=\fontsize{18.2pt}{23.7pt}\selectfont]
\draw  (5,9) circle (1.25cm);
\draw  (15,9) circle (1.25cm);
\draw [-{Stealth[scale=1.5]}, ] (6.25,10.25) -- (13.75,10.25);
\draw [-{Stealth[scale=1.5]}, ] (13.75,7.75) -- (6.25,7.75);
\node [font=\fontsize{18.2pt}{23.7pt}\selectfont, inner xsep=0.080cm, inner ysep=0.085cm, rounded corners=0.020cm] at (14.875,9) {UWB};
\node [font=\fontsize{18.2pt}{23.7pt}\selectfont, inner xsep=0.080cm, inner ysep=0.085cm, rounded corners=0.020cm] at (4.875,9) {UWB};
\node [font=\fontsize{18.2pt}{23.7pt}\selectfont, inner xsep=0.080cm, inner ysep=0.085cm, rounded corners=0.020cm] at (10,10.75) {POLL};
\node [font=\fontsize{15.9pt}{20.7pt}\selectfont, inner xsep=0.080cm, inner ysep=0.085cm, rounded corners=0.020cm] at (5.75,8.875) {1};
\node [font=\fontsize{15.9pt}{20.7pt}\selectfont, inner xsep=0.080cm, inner ysep=0.085cm, rounded corners=0.020cm] at (15.75,8.875) {2};
\node [font=\fontsize{18.2pt}{23.7pt}\selectfont, inner xsep=0.080cm, inner ysep=0.085cm, rounded corners=0.020cm] at (10,7.25) {RESP};
\node [font=\fontsize{18.2pt}{23.7pt}\selectfont, inner xsep=0.080cm, inner ysep=0.085cm, rounded corners=0.020cm] at (6.125,10.875) {$t_{x1}$};
\node [font=\fontsize{18.2pt}{23.7pt}\selectfont, inner xsep=0.080cm, inner ysep=0.085cm, rounded corners=0.020cm] at (6.125,7.25) {$r_{x1}$};
\node [font=\fontsize{18.2pt}{23.7pt}\selectfont, inner xsep=0.080cm, inner ysep=0.085cm, rounded corners=0.020cm] at (13.875,10.875) {$r_{x2}$};
\node [font=\fontsize{18.2pt}{23.7pt}\selectfont, inner xsep=0.080cm, inner ysep=0.085cm, rounded corners=0.020cm] at (13.875,7.25) {$t_{x2}$};
\end{circuitikz}
}%
\caption{TWR explanation diagram}
\label{fig:TWR}
\end{figure}
\noindent The algorithm works in 4 steps:  
\begin{enumerate}
    \item $UWB_1$ sends a \textbf{POLL} to $UWB_2$. The time of emission $t_{x1}$ is memorized by $UWB_1$.
    \item $UWB_2$ receives the \textbf{POLL} and memorizes the time of arrival $r_{x2}$.
    \item $UWB_2$ processes the \textbf{POLL} and precisely plans the time of emission $t_{x2}$ for the \textbf{RESP}. It contains ($t_{x2},r_{x2}$).
    \item $UWB_1$ receives the \textbf{RESP} and memorizes $r_{x1}$.
\end{enumerate}
With all these parameters, we can deduce the following variables:
\begin{itemize}
    \item Total time: $t_{tot} = r_{x1} - t_{x1}$
    \item Processing time of $UWB_2$: $t_{proc} = r_{x2} - t_{x2}$
    \item Travel time: $t_{travel} = t_{tot} - t_{proc}$
\end{itemize}
The signal travels back and forth, so it covers twice the distance ($d$), and the travel speed of a radio signal is equivalent to the speed of light ($c$). \\ 
We can therefore deduce this equation: 
\begin{equation} \label{TWR_eq}
\boxed{\frac{t_{travel}\times c}{2}=\frac{(t_{tot} - t_{proc}) \times c}{2}= \frac{((r_{x1} - t_{x1}) - (r_{x2} - t_{x2})) \times c}{2}}
\end{equation}
We can also determine the error induced by this algorithm.\\
After calculation (\ref{TWR-Error-calculation}), we obtain this equation: 
\begin{equation}
\varepsilon \approx \frac{t_ {proc}(e_a - e_b)}{2}
\end{equation}
As we can observe, the processing time is directly linked to the error. The $t_{proc}$ is around 3,000,000 nanoseconds (3 ms), so it induces a high error.

\subsubsection{Double Sided Two way ranging (DS-TWR)}
\label{DS-TWR}
In the TWR algorithm, the main drawback is the error and inaccuracy.
To remove that error, the DS-TWR can be implemented. This algorithm also uses 2 modules $UWB_1$ and $UWB_2$.

\begin{figure}[!ht]
\centering
\resizebox{0.7\textwidth}{!}{%
\begin{circuitikz}
\tikzstyle{every node}=[font=\fontsize{18.2pt}{23.7pt}\selectfont]
\draw  (5,9) circle (1.25cm);
\draw  (15,9) circle (1.25cm);
\draw [-{Stealth[scale=1.5]}, ] (13.75,7.75) -- (6.25,7.75);
\node [font=\fontsize{18.2pt}{23.7pt}\selectfont, inner xsep=0.080cm, inner ysep=0.085cm, rounded corners=0.020cm] at (14.875,9) {UWB};
\node [font=\fontsize{18.2pt}{23.7pt}\selectfont, inner xsep=0.080cm, inner ysep=0.085cm, rounded corners=0.020cm] at (4.875,9) {UWB};
\node [font=\fontsize{18.2pt}{23.7pt}\selectfont, inner xsep=0.080cm, inner ysep=0.085cm, rounded corners=0.020cm] at (10,9.5) {POLL};
\node [font=\fontsize{15.9pt}{20.7pt}\selectfont, inner xsep=0.080cm, inner ysep=0.085cm, rounded corners=0.020cm] at (5.75,8.875) {1};
\node [font=\fontsize{15.9pt}{20.7pt}\selectfont, inner xsep=0.080cm, inner ysep=0.085cm, rounded corners=0.020cm] at (15.75,8.875) {2};
\node [font=\fontsize{18.2pt}{23.7pt}\selectfont, inner xsep=0.080cm, inner ysep=0.085cm, rounded corners=0.020cm] at (10,7.25) {RESP};
\node [font=\fontsize{18.2pt}{23.7pt}\selectfont, inner xsep=0.080cm, inner ysep=0.085cm, rounded corners=0.020cm] at (6.125,10.875) {$t_{x1}$};
\node [font=\fontsize{18.2pt}{23.7pt}\selectfont, inner xsep=0.080cm, inner ysep=0.085cm, rounded corners=0.020cm] at (6.125,7.25) {$r_{x1}$};
\node [font=\fontsize{18.2pt}{23.7pt}\selectfont, inner xsep=0.080cm, inner ysep=0.085cm, rounded corners=0.020cm] at (13.875,10.875) {$r_{x2}$};
\node [font=\fontsize{18.2pt}{23.7pt}\selectfont, inner xsep=0.080cm, inner ysep=0.085cm, rounded corners=0.020cm] at (13.875,7.25) {$t_{x2}$};
\draw [{Stealth[scale=1.5]}-, ] (13.625,10.125) -- (6.125,10.125);
\draw [{Stealth[scale=1.5]}-, ] (13.5,11.75) -- (6,11.75);
\node [font=\fontsize{18.2pt}{23.7pt}\selectfont, inner xsep=0.080cm, inner ysep=0.085cm, rounded corners=0.020cm] at (9.625,11.25) {FINAL};
\node [font=\fontsize{18.2pt}{23.7pt}\selectfont, inner xsep=0.080cm, inner ysep=0.085cm, rounded corners=0.020cm] at (6.125,12.5) {$t_{x12}$};
\node [font=\fontsize{18.2pt}{23.7pt}\selectfont, inner xsep=0.080cm, inner ysep=0.085cm, rounded corners=0.020cm] at (13.75,12.5) {$r_{x22}$};
\end{circuitikz}
}%
\caption{DS-TWR explanation diagram}
\label{fig:DS-TWR}
\end{figure}

\noindent The algorithm works in 6 steps :  
\begin{enumerate}
    \item $UWB_1$ sends a \textbf{POLL} to $UWB_2$. The time of emission
          $t_{x1}$ is memorized by $UWB_1$.
    \item $UWB_2$ receives the \textbf{POLL} and memorizes the time of
          arrival $r_{x2}$.
    \item $UWB_2$ processes the \textbf{POLL} and precisely plans the time
          of emission $t_{x2}$ of the \textbf{RESP}, so that the reply
          delay $D_b = t_{x2} - r_{x2}$ is known.
    \item $UWB_1$ receives the \textbf{RESP} and memorizes the time of
          arrival $r_{x1}$.
    \item $UWB_1$ processes the \textbf{RESP} and precisely plans the time
          of emission $t_{x12}$ of the \textbf{FINAL}, so that the reply
          delay $D_a = t_{x12} - r_{x1}$ is known. The \textbf{FINAL}
          contains $(t_{x1},\, r_{x1},\, t_{x12})$.
    \item $UWB_2$ receives the \textbf{FINAL} and memorizes the time of
          arrival $r_{x22}$.
\end{enumerate}

\noindent With all those parameters, we can compute four time intervals :
\begin{itemize}
    \item Response time for $UWB_1 : t_1 = r_{x1} - t_{x1}$
    \item Response time for $UWB_2 : t_2 = r_{x22} - t_{x2}$
    \item Processing time of $UWB_1 : t_{proc_1} = t_{x12} - r_{x1}$
    \item Processing time of $UWB_2 : t_{proc_2} = r_{x2} - t_{x2}$
\end{itemize}

\noindent We can deduce the total travel time: 

\begin{equation}
\boxed{t_{travel} = \frac{t_1 \times t_2 - t_{proc_1} \times t_{proc_2}}{ t_1 +t_2+t_{proc_1}+t_{proc_2}}}
\end{equation}

\noindent And then find the distance between the 2 modules : 
 \begin{equation}
    d = c \times t_{travel}
\end{equation}         

\noindent Now let's compare the induced error. After calculation (\ref{DSTWR-Error-calculation}), we obtain this equation  : 
\begin{equation}
\varepsilon \approx T_f \times \frac{e_a + e_b}{2}
\end{equation}
As we can see, the error is not multiplied by $t_{proc}$, so the error is reduced.\\
With this algorithm, we are not dependent on the processor anymore. That is the algorithm currently used in our project.

\subsection{Timing algorithm}\label{UWB_Algo_timings}
When working on a drone swarm, there are not only two UWB modules : there are actually multiple modules trying to communicate. One weakness of the UWB is its sensitivity to noise and interference. The "way" should be clear before emitting, otherwise that creates collisions and damages the signal, ultimately leading to corrupted distance estimations. For this purpose, we need to look at how to divide the time of emission to avoid collision. In our specific case, we are working with 3 to 4 drones.
\subsubsection{ALOHA}
This method is named after the first ALOHAnet system, which introduced it in 1970. It provides a simple approach to channel access control without requiring any time synchronization. Instead of planning transmissions, a module sends its data packet immediately, regardless of whether the medium is currently busy or free.
The primary advantage of  ALOHA lies in its ease of implementation and full decentralization, allowing independent nodes or drones to join the network dynamically without prior initialization. 
However, its major drawback is its susceptibility to collisions caused by uncoordinated simultaneous transmissions. As demonstrated in foundational studies, the maximum theoretical throughput (channel capacity) of  ALOHA is limited to approximately $18.4\%$ \cite{noauthor_pdf_nodate}. This low efficiency makes it unsuitable for high-frequency localization and critical communication applications.


\subsubsection{Time-division multiple access (TDMA)} \label{TDMA}
This method aims to to divide a fixed period of time in frames, and cut it in different slots. Each slot is dedicated to each drone. Those frames and slots are fixed in the hardware and are stricly respected. 
\begin{figure}[H]
    \centering
    \resizebox{0.7\textwidth}{!}{
    \begin{circuitikz}
    \draw  (0.5,12.875) rectangle (2.5,11.25);
    \draw  (2.5,12.875) rectangle (4.5,11.25);
    \draw  (4.5,12.875) rectangle (6.5,11.25);
    \draw  (6.5,12.875) rectangle (8.5,11.25);
    \draw  (8.5,12.875) rectangle (10.5,11.25);
    \draw [short] (0.5,13.25) -- (10.5,13.25);
    \draw [short] (0.25,13) arc[start angle=178.78, end angle=91.22, radius=0.2555];
    \draw [short] (10.5,13.25) arc[start angle=90.00, end angle=0.00, radius=0.2500];
    \draw [short] (10.625,13.625) -- (10.625,13.625);
    \draw [short] (10.625,13.625) -- (10.625,13.625);
    \draw [short] (5.75,14) -- (5.75,14);
    \draw [short] (5.75,14) -- (5.75,14);
    \draw [short] (5.875,13.875) -- (5.875,13.875);
    \node [font=\fontsize{18.2pt}{23.7pt}\selectfont, inner xsep=0.080cm, inner ysep=0.085cm, rounded corners=0.020cm] at (5.5,13.625) {Frame};
    \node [font=\fontsize{18.2pt}{23.7pt}\selectfont, inner xsep=0.080cm, inner ysep=0.085cm, rounded corners=0.020cm] at (1.5,12) {T1};
    \node [font=\fontsize{18.2pt}{23.7pt}\selectfont, inner xsep=0.080cm, inner ysep=0.085cm, rounded corners=0.020cm] at (3.5,12) {T2};
    \node [font=\fontsize{18.2pt}{23.7pt}\selectfont, inner xsep=0.080cm, inner ysep=0.085cm, rounded corners=0.020cm] at (5.5,12) {T3};
    \node [font=\fontsize{18.2pt}{23.7pt}\selectfont, inner xsep=0.080cm, inner ysep=0.085cm, rounded corners=0.020cm] at (7.5,12) {T4};
    \node [font=\fontsize{18.2pt}{23.7pt}\selectfont, inner xsep=0.080cm, inner ysep=0.085cm, rounded corners=0.020cm] at (9.5,12) {T5};
    \end{circuitikz}
    }
    \caption{TDMA explanation diagram}
\end{figure}

For example, if we have 3 drones, the frame of $UWB_1$, when it is looking for distances with a DS-TWR algorithm, is divided in 5 parts.
\begin{enumerate}
    \item T1: $UWB_1$ sends the \textbf{POLL}  to $UWB_2$ and  $UWB_3$
    \item T2: $UWB_2$ sends the \textbf{RESP}  to $UWB_1$
    \item T3: $UWB_1$ sends the \textbf{FINAL} to  $UWB_2$
    \item T4: $UWB_3$ sends the \textbf{RESP}  to $UWB_1$
    \item T5: $UWB_1$ sends the \textbf{FINAL} to $UWB_3$
\end{enumerate}

\noindent This method has a clear advantage: there is no collision at all due to strict time periods. It also permits to shut down the modules when we do not need them and save energy. However, there are a few drawbacks: firstly, if we add one more drone, we have to add two other slots, and when reaching 10 drones, there is no time anymore to get information at a high frequency. Moreover, it needs a synchronization to make all the modules' clocks work at the same speed and avoid collision perfectly.

\subsubsection{Slotted ALOHA}

This method combines the features of both the ALOHA and TDMA protocols. Time is divided into slots, and UWB modules are only allowed to start their transmissions at the beginning of a slot boundary. If a module has data ready to transmit, but a slot has already been taken, it must wait for the next available slot.

This approach preserves the flexibility of ALOHA while doubling the maximum theoretical throughput to approximately $36.8\%$ \cite{noauthor_pdf_2026}. However, it still remains less deterministic and efficient than TDMA due to potential collisions when multiple nodes choose the same slot.
\newpage
\section{Measure processing}

In most cases, the data measured by sensors are noisy. It could come from a sensor offset, or from crosstalk or other exterior interactions. In order to obtain clean measurements and to be able to use them for controlling or any other purpose, we have to develop  filters.

\subsection{Extended Kalman Filter for relative localisation  }
In the context of navigating systems in noisy environments with sensors, the Kalman filter is a relevant algorithm. As Welch and Bishop (2006, p. 1) [3] point out, the Kalman filter is notable for its ability to "support estimations of past, present, and even future states [...] even when the precise nature of the modeled system is unknown.". In our study case, we deal with a drone swarm communicating using the UWB communication protocol, which gives us the relative distance between each pair of drones by using the time of flight of a radio wave. The UWB therefore provides a scalar distance, whereas the quantity of interest is the relative position vector. The measurement model, which links the two through the Euclidean norm, is nonlinear, which calls for the extended Kalman filter. This version resorts to a typical technique of linearization around the current state using a Jacobian matrix. We are also considering an anchor-free system which enhances the challenge of this project. In this section, we will first cover the theory of the classical Kalman filter and its mathematical formalist and then dive into the extended Kalman filter and its utility.


\subsubsection{Theory of the Kalman filter}

\paragraph{Origin of the theory}

In his pioneering article of 1960 \parencite{kalman1960}, 
Rudolf Kalman describes state estimation as a recursive process in which at each time step, 
the current estimation and the uncertainty are computed using only the past estimate and the newly available measure. When assuming a linear dynamical system and Gaussian noise, the Kalman filter minimizes the mean-square estimation error : in other words, the best possible estimate given the available measurements. Even when these assumptions do not hold exactly, the filter remains the best linear estimator achievable (known as the Best Linear Unbiased Estimator, or BLUE). This optimality is what established the Kalman filter as the standard approach for state estimation in noisy dynamic systems. While Kalman filter's optimality holds for linear dynamical systems, we have in our case, as mentioned earlier, a non-linearity with the use of the Euclidian distance. This is why we require an extended version of the filter that can handle non-linearity, even at the cost of losing strict optimality.


\paragraph{Formalism of the discrete Kalman Filter}

Let's note $x_k\in \mathbb{R}_n$ the current state of a discrete system governed by the following linear equation :
$${\displaystyle \mathbf {x} _{k}=\mathbf {F} _{k}\mathbf {x} _{k-1}+\mathbf {B} _{k}\mathbf {u} _{k}+\mathbf {w} _{k}}$$

In which we have :

\begin{itemize}
    \item $F_k$ is the state transition matrix, linking the previous state $x_{k-1}$ to the current state $x_k$
    \item $B_k$ is the control-input matrix, linking the control vector $u_k$ to the current state $x_k$
    \item $w_k$ is the process noise which is assumed to be Gaussian and centered, with a covariance matrix $Q_k$
\end{itemize}

At each time $k$, the measurement $z_k$ is written as follows :
$${\displaystyle \mathbf {z} _{k}=\mathbf {H} _{k}\mathbf {x} _{k}+\mathbf {v} _{k}}$$
\newpage
With :
\begin{itemize}
    \item $H_k$ is the observation matrix
    \item $w_k$ is the observation noise which is assumed to be Gaussian and centered, with a covariance matrix $R_k$
\end{itemize}

We can divide the Kalman filter process into 2 phases : the "prediction" phase and the "update" phase. 
The first one does an estimate of the state at the next time step by using the estimate of 
the current time step and the error covariance. The second phase can be thought of as a feedback : it 
includes a new measurement in the current state to improve its estimate.

\begin{figure}[H]
    \centering
    \includegraphics[width=0.6\linewidth]{kalman_2_phases.png}
    \caption{Ongoing discrete Kalman filter cycle  (Welch and Bishop (2006, p. 5) \parencite{welch2006}) }
    \label{fig:placeholder}
\end{figure}
\subsubsection{Extended Kalman filter}


As mentioned earlier, the extended Kalman filter is used to deal with a non-linear state process defined as diffentiable function $f$. To do so, it essentially linearizes the nonlinear function around the current estimate by computing the Jacobian matrices of the nonlinear functions $f$ and $h$ in the following equations (with the same notations as before) :

$$ {\displaystyle {\textbf {x}}_{k}=f({\textbf {x}}_{k-1},{\textbf {u}}_{k},{\textbf {w}}_{k})}$$
$${\displaystyle {\textbf {z}}_{k}=h({\textbf {x}}_{k},{\textbf {v}}_{k})}$$


Like the original filter, it is recursive: only the current estimate and its covariance are stored, which keeps memory and computation low enough for embedded platforms, while the covariance provides a measure of uncertainty that can be used to reject inconsistent measurements. However, this comes with a drawback. The linearization is a first-order approximation, so the optimality of the linear Gaussian case is lost, and the filter may diverge if the initial estimate is poor or the system is strongly nonlinear \parencite{boutayeb_convergence_1997}. Its performance also depends on the tuning of the noise covariances. Alternatives such as the unscented Kalman filter \parencite{julier_new_1997}, particle filters \parencite{arulampalam_tutorial_2002} or optimization-based estimators differ in accuracy, robustness and computational cost. The EKF nevertheless remains the standard approach in navigation, including in open-source flight controllers (like in the ArduPilot firmware).
\newpage

\subsection{UWB localization : benefits and limits }

As described in the previous section, UWB ranging enables centimeter-level accuracy via time of flight calculations. However, the data gives us only the relative distances between each drone which is not enough to perform localisation : in particular, information about orientation is missing. In this section, we will discuss the viability of a distance-only approach and the topic of observability.

\subsubsection{Range-only system}

In their paper entitled "Dynamic Range-Only Localization for Multi-Robot Systems", Cao et al. (2018) \parencite{cao_2018} propose a method to perform robot localization via range measurement only. To do so, they present a method based on successive stages which allows to dynamically localize multiple robots by using trilateration and an EKF recursive estimation. While this article shows a functional multi-agent system with relative localization without anchors, it has been demonstrated only in simulation and in a 2D environment. However, it does not address experimental validation and three-dimensional flight. Moreover, their approach uses a coordinate system built by the robots themselves while drones are usually localized in an Earth-fixed frame via their flight controllers. The authors acknowledge the question of a global frame but consider it outside the scope of their work, mentioning that a swarm deployed in an unknown environment might no have access to global positioning.

\subsubsection{Observability}

As defined in \parencite{astrom_2021}, we say that "a system is \textit{observable} if for every $T>0$, it is possible to determine the state of the system $x(T)$ through measurements of the $y(t)$ and $u(t)$ in the interval $[0,T]$" where $y(t)$ and $u(t)$ are respectively the output and the input of the process at time $t$. Observability plays an important role in control theory because it is common that is all the states are not accessible through measurement : it is therefore possible to build an \textit{observer}, a dynamical system represented by a mathematical model capable of recreating the states of the system. 

In our project, the observability problem is harder to deal with because it would be necessary to know which states of the drones we choose to consider. Van der helm et al. \parencite{van_der_helm_2020} studied a similar problem : they implemented UltraWide band technology for both ranging and communication between Parrot drones (velocities, accelerations, yaw rates and height). They tested different configurations of leader-follower flights and performed observability analyses. They showed, on the one hand, that without a common heading, "the motion
of agents must meet more stringent conditions to be observable and the relative localization becomes more susceptible to noise on the range measurements". On the other hand, the dependency on a common heading offers the advantage of being less sensitive to local disturbances in Earth's magnetic field, which, even when small, can lead to a large negative impact on the EKF used. In a leader-follower configuration, a fixed-geometry (like 2 drones moving in parallel) is not possible and for that, they implement a method where the follower flies in a time-delayed trajectory version of the leader. Their approach was validated indoors with an on-board EKF but outdoor operations with GNSS are not addressed.
\subsection{Data fusion}

\subsubsection{Data fusion with an EKF in drone swarm applications}
In the context of our project, the most relevant article is most likely "Onboard Ranging-Based Relative Localization and Stability for Lightweight Aerial Swarms" by Li et al. \parencite{li_2025} from TU Delft in the Netherlands.
\newpage
In their paper, they share an autonomous relative localization technique fusing UWB distance measurements and state information like velocity, yaw rate, height from neighbours. The following illustration shows their scheme of work : 

\begin{figure}[H]
    \centering
    \includegraphics[width=0.65\linewidth]{Li_2025.png}
    \caption{Scheme of work of Li et al. \parencite{li_2025} }
    \label{fig:placeholder}
\end{figure}

By using the data of 3 sensors (IMU, optical flow, laser), they retrieve the 3-axis velocity $\overline{v_i}$, the pitch $\theta$ and the roll $\phi$ in the body frame of the drone. The yaw rate $\overline{r}$, the range $d$ (distance between 2 drones) and the height $h$ are respectively given by the gyroscope, the UWB sensors and the laser. One of the main aspects of their project is that they work with tiny, lightweight (33 grams) drones which are resource-constrained. Therefore, an Extended Kalman filter is exploited to combine the data, particularly because it has a low computational overhead. The article is particularly appropriate because it deals with every aspect of our project. It brings up the fusion of multiple sensors through an EFK, the use of Ultra-wide band for both ranging and communication as well as observability considerations and hardware resource constraints. Nevertheless, they do not implement GPS when merging data, which differs from the context of our work. 

\subsubsection{GPS consideration}

The papers mentioned previously did not include GPS localization in their architecture but our drones do include a GPS tracker, which we would like to take advantage of. We are working on an outdoor system and it is possible through the flight controller to use the GPS data to get a more accurate state estimate of the drone. In an article of 2020, Chen et al. \parencite{chen_2020} studied the complimentarity of UWB measurements with GPS localization to get a better accuracy. They show a great improvement of device localization with an accuracy up to 10 centimeters. This approach is different from ours as the GPS data is directly combined through a cost function with the UWB measurements, while we are planning to use an Extended Kalman filter. Additionally, their experiments were made on static objects and the final localization is obtained in a 2D plane, while we deal with moving UAVs in a 3D frame, which may add more challenges. 
\\

A closer reasearch to our context of study was made by Yang et Lee in 2023 \parencite{yang_2023} in which they implemented an Extended Kalman Filter using GNSS and INS (Inertial Navigation System) with complementary UWB measurements in a flight controller firmware. They used the flight controller "Pixhawk 4" (the same as us)  and successfully carried out a field test which, however, did not include a moving drone formation. 

\section{Drone Swarm Formation Control}

After considering the data processing, it is necessary to understand how we can leverage it to obtain a functional drone swarm formation control.
This topic is very broad and multiple methods of control can be considered : for the sake of simplicity, we will try to cover methods relevant for our applications which do not require to dive too much into mathematical considerations.
We will first study the different possible topologies of swarms and then we will focus on formation control strategies. 


\subsection{Formation Control Scheme}

We can firstly categorize multiple formation control schemes depending on where the decisions are made and how information is shared in the swarm. In their article of 2024 \parencite{bu_2024}, Bu et al. did an exhaustive review of UAV swarm formation control and defined three main formation control schemes : centralized, decentralized and distributed. In a centralized scheme, a single control point such as a computer is in charge of managing and processing all the information as well as taking decisions and giving commands to individual UAVs. In this configuration, the drones' autonomy is thus limited. While this configuration can benefit from simpler decision-making and easier implementation, it suffers from scalability issues and a single point of failure problem. In a decentralized approach, each drone acts based on its local perception and predefined rules and usually limits inter-drone communication. In this case, formation coordination and design are more challenging but scalability and adaptability are improved. Finally, the distributed scheme resorts to a more collaborative approach and hence, an extensive communication between the UAVs for teamwork. This configuration is relevant in applications where the use of collective information from different sensors is primordial.



\subsection{Formation Control Strategy}

Control strategies include the definition of the roles within a swarm as well as the way the swarm is controlled. This broad topic encompasses a lot of aspects that we will not cover exhaustively in this study, particularly on the mathematical theory behind the control implementation.


In the review of Bu et al. \parencite{bu_2024} mentioned previously, they discuss various conventional strategies and their characteristics. An intuitive way of imagining the configuration of a swarm involves giving each drone a role of either leader or follower. The basic idea is to have one or multiple designated leader drones following the given mission while follower drones maintain a position or a distance from the leaders or other followers. This method shows a good flexibility and can be applied in many different contexts. However, it heavily relies on the leader which could cause heavy problems if it is affected. Other approaches include virtual-structures, behaviour-based architectures, consensus-based architectures and artificial potential fields. Within the context of our project, our main constraint lies on the hardware, where the implemented control strategy should run on a ESP32 microcontroller. The leader-follower architecture reduces the overall computational cost and complexity of the system since the leader does not necessarily need to be part of the control architecture.
In order to keep a pre-defined formation, each drone within the swarm needs to constantly adjust itself to be in the right position. As defined in the article of Oh et al. \cite{oh_survey_2015}, it is possible to categorize the types of formation control depending on the sensing capabilities and the interaction topology of agents.
\newpage
The categories can be represented as follows :


\begin{figure}[H]
        \centering
        \includegraphics[width=0.6\linewidth]{formation_control_Oh.png}
        \caption{Sensing capability vs. interaction topology \parencite{oh_survey_2015}}
        \label{fig:placeholder}
\end{figure}
       
Each control architecture can be summarized as follows : 

\begin{itemize}
    \item Position-based control: Agents control their own positions within a global coordinate system to achieve the desired formation.
    \item Displacement-based control: Agents regulate relative displacements toward their neighbours using a shared global orientation, without needing to know their own absolute coordinates.
    \item Distance-based control: Agents regulate purely inter-agent distances using individual local coordinate systems that do not need to be aligned with one another.
\end{itemize}

Our project involves multiple sensors such as a GNSS module, a magnetometer as well as distance ranging  and communication abilities. Theoretically, all the types of formation control could be implemented, however it is important to consider other factors. The most relevant and robust algorithm for our system is the distance-based control as it does not need any sort of global coordinate system or orientation to work and can rely exclusively on relative data. Moreover, our project takes advantage of the UWB sensors for distance measurements. It is therefore more interesting to put it to the test by basing the formation control on it. 

\newpage
\section{Communication protocol}
 
\subsection{Xbee}

XBee is a family of radio modules made by Digi International. Depending on the model and firmware, it implements IEEE 802.15.4 \cite{noauthor_ieee_nodate} directly, ZigBee, or Digi’s own DigiMesh protocol. The module is controlled by a microcontroller through a simple serial link (UART or SPI), either in transparent mode (the module behaves like a wireless serial cable) or in API mode (the host sends and receives structured frames with addresses and acknowledgements). \\

\noindent Two bands are commonly used. The 2.4 GHz modules have a high data rate and are widely available. The 900 MHz modules trade throughput for a much longer range. At the MAC layer, 802.15.4 \cite{noauthor_ieee_nodate} uses CSMA/CA, so a module listens before transmitting, and unicast frames can be acknowledged and retried automatically.

\begin{table}[H]
\centering

\label{tab:xbee}
\begin{tabular}{@{}p{3.7cm}p{10cm}@{}}
\toprule
\textbf{Criterion} & \textbf{XBee} \\
\midrule
Range & $\sim$ 60\,m indoor, $\sim$1\,km outdoor line-of-sight  \\
Latency & A few ms to tens of ms per hop  \\
Data rate & 250\,kbps raw  \\
Power consumption & $\sim$ 30--45\,mA in RX/TX (standard module)\\
Topology & Point-to-point, star, mesh (ZigBee, DigiMesh) \\
Number of nodes & Hundreds (16-bit addressing)\\

\bottomrule
\end{tabular}
\caption{XBee characteristics}
\end{table}


\subsection{UWB}
A UWB module is a radio transceiver, so the frames used for ranging (POLL, RESP, FINAL) can carry a small payload. In the DS-TWR algorithm described above, the FINAL frame already contains timestamps (tx1, rx1, tx12). The same mechanism can carry a drone ID, a status flag or a short command. The advantage is that localization and data exchange share a single radio and a single time schedule, so no second module is needed for small messages.\\

\noindent The main limitation is the frame size and the energy cost. UWB frames are short (about 127 bytes in the standard configuration), and the receiver is power-hungry compared with the other technologies in this section. UWB is therefore well suited to short control messages, but not to image or bulk data transfer.

\begin{table}[H]
\centering
\label{tab:uwb}
\begin{tabular}{@{}p{3.7cm}p{10cm}@{}}
\toprule
\textbf{Criterion} & \textbf{UWB} \\
\midrule
Range & $\sim$10--50\,m in practice indoors, up to $\sim$300\,m line-of-sight \\
Latency & Frame airtime under 1\,ms; a full DS-TWR exchange takes a few ms \\
Data rate & 110\,kbps, 850\,kbps or 6.8\,Mbps; small payload per frame ($\sim$127 bytes) \\
Power consumption & $\sim$ 30--40\,mA in TX; RX is higher \\
Topology & Peer-to-peer / star (tags and anchors); no native mesh \\
Number of nodes & Limited by the channel access method (ALOHA, TDMA); a few tens with ranging \\
\bottomrule
\end{tabular}\caption{UWB characteristics }
\end{table}

\subsection{ESP-NOW}
ESP-NOW is a protocol developed by Espressif for its ESP32 and ESP8266 chips. It uses Wi-Fi radio hardware at 2.4 GHz, but without a router, an access point or an IP connection: data is sent in short vendor-specific Wi-Fi frames directly from one device to another, identified by their MAC addresses. A device registers its peers, sends a packet (up to 250 bytes in the classic version), and receives a callback telling it whether the frame was acknowledged. Broadcast is also possible, and encryption can be enabled per peer. The ESP-NOW connection is connectionless, so there is no handshake before sending. This explains the low latency.\\

\noindent Studies agree on its main features \cite{eridani_comparative_2021}. Authors compared ESP-NOW with Wi-Fi and Bluetooth on range, speed, latency, energy usage and barrier resistance. They also evaluated the maximum distance and power consumption of ESP-NOW nodes. A community test \cite{noauthor_esp-now_nodate} published by Espressif reports latency under 20 ms in open field up to 300 m, and a long-range mode (ESP-NOW-LR) reaching up to about a kilometre. Note that this source is not peer-reviewed. 

\begin{table}[H]
\centering
\label{tab:espnow}
\begin{tabular}{@{}p{3.7cm}p{12cm}@{}}
\toprule
\textbf{Criterion} & \textbf{ESP-NOW} \\
\midrule
Range & $\sim$ 15--50\,m indoors; $\sim$ 200--300\,m outdoors line-of-sight; up to $\sim$1\,km  \\
Latency & Typically a few ms; under $\sim$ 20\,ms measured in open field \\
Data rate & Depends on the configured Wi-Fi rate (1\,Mbps by default), 250 bytes per packet \\
Power consumption & $\sim$ 100--240\,mA in RX/TX;\\
Topology & Peer-to-peer, one-to-many, broadcast \\
Number of nodes & About 20 registered peers per device; broadcast has no such limit \\
\bottomrule
\end{tabular}\caption{ESP-NOW characteristics}
\end{table}

\newpage
